% 卡斯蒂隆问题（Castillon's Problem）
% 已知一个圆及不在圆上的三点，需要求作这个圆的内接三角形，
% 使得它的三边或其延长线分别通过这三个已知点。

\documentclass{standalone}
\usepackage{tikz}

\usetikzlibrary{cartes}

\makeatletter
\tikzset{
  % 构造变换: 从圆上的一点 P，通过给定点 A 射影到圆上另一点 P'
  % 再通过给定点 B 射影到圆上另一点P''，通过给定点C射影到圆上另一点P'''
  % 设从圆上点 P 到圆上另一个点 P''' 的变换为这 3 次射影的合成
  % 问题转换为寻找该变换的不动点
  % single transform={\Circle,\A,\P,\Q},
  single transform/.code args={#1,#2,#3,#4}{
   \tikzset{
      conics/intersect cl={#1,#2,#3,\ca@U,\ca@V},
      conics/exclude={\ca@U,\ca@V,#3,#4},
    }
  },
  % triple transform={\Circle,\A,\B,\C,\P,\Q}
  triple transform/.code args={#1,#2,#3,#4,#5,#6}{
   \tikzset{
      single transform={#1,#2,#5,\ca@S},
      single transform={#1,#3,\ca@S,\ca@T},
      single transform={#1,#4,\ca@T,#6},
    }
  },
}
\makeatother

\tikzmath{
   % 点采用齐次坐标
   \O{1,1} = 0; \O{2,1} = 0; \O{3,1} = 1;
   \radius = 2;
   \A{1,1} = 1; \A{2,1} = 1; \A{3,1} = 1;
   \B{1,1} = -0.5; \B{2,1} = 0.5; \B{3,1} = 1;
   \C{1,1} = -3; \C{2,1} = -2; \C{3,1} = 1;
   \Pa{1,1} = \O{1,1} + \radius * cos(-90);
   \Pa{2,1} = \O{2,1} + \radius * sin(-90);
   \Pa{3,1} = 1;
   \Qa{1,1} = \O{1,1} + \radius * cos(-45);
   \Qa{2,1} = \O{2,1} + \radius * sin(-45);
   \Qa{3,1} = 1;
   \Ra{1,1} = \O{1,1} + \radius * cos(-60);
   \Ra{2,1} = \O{2,1} + \radius * sin(-60);
   \Ra{3,1} = 1;
}

\tikzset{
  conics/define/circle/center-radius={\O,\radius,\Circle},
  conics/dehomogenize={\A,\AA},
  conics/dehomogenize={\B,\BB},
  conics/dehomogenize={\C,\CC},
}

\begin{document}

\begin{tikzpicture}
  \tikzset{
    single transform={\Circle,\A,\Pa,\Pb},
    single transform={\Circle,\B,\Pb,\Pc},
    single transform={\Circle,\C,\Pc,\Pd},
    triple transform={\Circle,\A,\B,\C,\Pa,\U},
    conics/dehomogenize={\Pa,\PPa},
    conics/dehomogenize={\Pb,\PPb},
    conics/dehomogenize={\Pc,\PPc},
    conics/dehomogenize={\Pd,\PPd},
    conics/dehomogenize={\U,\UU},
  }
  \axes[grid] (-4:4,-4:4);
  \coordinate (A) at (\AA{1},\AA{2});
  \coordinate (B) at (\BB{1},\BB{2});
  \coordinate (C) at (\CC{1},\CC{2});
  \coordinate (P1) at (\PPa{1},\PPa{2});
  \coordinate (P2) at (\PPb{1},\PPb{2});
  \coordinate (P3) at (\PPc{1},\PPc{2});
  \coordinate (P4) at (\PPd{1},\PPd{2});
  \coordinate (U) at (\UU{1}, \UU{2});

  \conic[draw=navy,thick] (\Circle);
  \segment[draw=blue,clip=(-4:4,-4:4)] (\Pa,\A);
  \segment[draw=blue,clip=(-4:4,-4:4)] (\Pb,\B);
  \segment[draw=blue,clip=(-4:4,-4:4)] (\Pc,\C);

  \foreach \p/\placement in {A/left,B/right,C/above,
  P1/above,P2/above,P3/above,P4/above,U/above}{
   \fill[red] (\p) circle (1.5pt);
   \draw (\p) node[\placement] {$\p$};
  }
\end{tikzpicture}

\begin{tikzpicture}
  \tikzset{
    single transform={\Circle,\A,\Qa,\Qb},
    single transform={\Circle,\B,\Qb,\Qc},
    single transform={\Circle,\C,\Qc,\Qd},
    conics/dehomogenize={\Qa,\QQa},
    conics/dehomogenize={\Qb,\QQb},
    conics/dehomogenize={\Qc,\QQc},
    conics/dehomogenize={\Qd,\QQd},
  }
  \axes[grid] (-4:4,-4:4);
  \coordinate (A) at (\AA{1},\AA{2});
  \coordinate (B) at (\BB{1},\BB{2});
  \coordinate (C) at (\CC{1},\CC{2});
  \coordinate (Q1) at (\QQa{1},\QQa{2});
  \coordinate (Q2) at (\QQb{1},\QQb{2});
  \coordinate (Q3) at (\QQc{1},\QQc{2});
  \coordinate (Q4) at (\QQd{1},\QQd{2});

  \conic[draw=navy,thick] (\Circle);
  \segment[draw=blue,clip=(-4:4,-4:4)] (\Qa,\A);
  \segment[draw=blue,clip=(-4:4,-4:4)] (\Qb,\B);
  \segment[draw=blue,clip=(-4:4,-4:4)] (\Qc,\C);

  \foreach \p/\placement in {A/left,B/right,C/above,
  Q1/above,Q2/above,Q3/above,Q4/above}{
   \fill[red] (\p) circle (1.5pt);
   \draw (\p) node[\placement] {$\p$};
  }
\end{tikzpicture}

\begin{tikzpicture}
  \tikzset{
    single transform={\Circle,\A,\Ra,\Rb},
    single transform={\Circle,\B,\Rb,\Rc},
    single transform={\Circle,\C,\Rc,\Rd},
    conics/dehomogenize={\Ra,\RRa},
    conics/dehomogenize={\Rb,\RRb},
    conics/dehomogenize={\Rc,\RRc},
    conics/dehomogenize={\Rd,\RRd},
  }
  \axes[grid] (-4:4,-4:4);
  \coordinate (A) at (\AA{1},\AA{2});
  \coordinate (B) at (\BB{1},\BB{2});
  \coordinate (C) at (\CC{1},\CC{2});
  \coordinate (R1) at (\RRa{1},\RRa{2});
  \coordinate (R2) at (\RRb{1},\RRb{2});
  \coordinate (R3) at (\RRc{1},\RRc{2});
  \coordinate (R4) at (\RRd{1},\RRd{2});

  \conic[draw=navy,thick] (\Circle);
  \segment[draw=blue,clip=(-4:4,-4:4)] (\Ra,\A);
  \segment[draw=blue,clip=(-4:4,-4:4)] (\Rb,\B);
  \segment[draw=blue,clip=(-4:4,-4:4)] (\Rc,\C);

  \foreach \p/\placement in {A/left,B/right,C/above,
  R1/above,R2/above,R3/above,R4/above}{
   \fill[red] (\p) circle (1.5pt);
   \draw (\p) node[\placement] {$\p$};
  }
\end{tikzpicture}

\begin{tikzpicture}
  \tikzset{
    triple transform={\Circle,\A,\B,\C,\Pa,\Pb},
    triple transform={\Circle,\A,\B,\C,\Qa,\Qb},
    triple transform={\Circle,\A,\B,\C,\Ra,\Rb},
    conics/cross={\Pa,\Qb,\la},
    conics/cross={\Pb,\Qa,\ma},
    conics/cross={\la,\ma,\M},
    conics/cross={\Pa,\Rb,\lb},
    conics/cross={\Pb,\Ra,\mb},
    conics/cross={\lb,\mb,\N},
    conics/cross={\Qa,\Rb,\lc},
    conics/cross={\Qb,\Ra,\mc},
    conics/intersect cl={\Circle,\M,\N,\Da,\Db},
    single transform={\Circle,\A,\Da,\Ea},
    single transform={\Circle,\B,\Ea,\Fa},
    single transform={\Circle,\A,\Db,\Eb},
    single transform={\Circle,\B,\Eb,\Fb},
    conics/dehomogenize={\A,\AA},
    conics/dehomogenize={\B,\BB},
    conics/dehomogenize={\C,\CC},
    conics/dehomogenize={\Da,\DDa},
    conics/dehomogenize={\Ea,\EEa},
    conics/dehomogenize={\Fa,\FFa},
    conics/dehomogenize={\Db,\DDb},
    conics/dehomogenize={\Eb,\EEb},
    conics/dehomogenize={\Fb,\FFb},
    conics/dehomogenize={\Pa,\PPa},
    conics/dehomogenize={\Pb,\PPb},
    conics/dehomogenize={\Qa,\QQa},
    conics/dehomogenize={\Qb,\QQb},
    conics/dehomogenize={\Ra,\RRa},
    conics/dehomogenize={\Rb,\RRb},
  }

  \axes[grid] (-4:4,-4:4);
  \coordinate (A) at (\AA{1},\AA{2});
  \coordinate (B) at (\BB{1},\BB{2});
  \coordinate (C) at (\CC{1},\CC{2});
  \coordinate (D1) at (\DDa{1},\DDa{2});
  \coordinate (E1) at (\EEa{1},\EEa{2});
  \coordinate (F1) at (\FFa{1},\FFa{2});
  \coordinate (D2) at (\DDb{1},\DDb{2});
  \coordinate (E2) at (\EEb{1},\EEb{2});
  \coordinate (F2) at (\FFb{1},\FFb{2});
  \coordinate (P1) at (\PPa{1},\PPa{2});
  \coordinate (P2) at (\PPb{1},\PPb{2});
  \coordinate (Q1) at (\QQa{1},\QQa{2});
  \coordinate (Q2) at (\QQb{1},\QQb{2});
  \coordinate (R1) at (\RRa{1},\RRa{2});
  \coordinate (R2) at (\RRb{1},\RRb{2});

  \conic[draw=navy,thick] (\Circle);
  \draw[thick,red] (D1) -- (E1) -- (F1) -- cycle;
  \draw[thick,violet] (D2) -- (E2) -- (F2) -- cycle;
  % \segment[draw=blue,clip=(-4:4,-4:4)] (\D,\A);
  % \segment[draw=blue,clip=(-4:4,-4:4)] (\D,\B);
  % \segment[draw=blue,clip=(-4:4,-4:4)] (\D,\C);
  % \segment[draw=teal,clip=(-4:4,-4:4)] (\E,\A);
  % \segment[draw=teal,clip=(-4:4,-4:4)] (\E,\B);
  % \segment[draw=teal,clip=(-4:4,-4:4)] (\E,\C);
  \segment[draw=teal,clip=(-4:4,-4:4)] (\la);
  \segment[draw=teal,clip=(-4:4,-4:4)] (\ma);
  \segment[draw=purple,clip=(-4:4,-4:4)] (\lb);
  \segment[draw=purple,clip=(-4:4,-4:4)] (\mb);
  \segment[draw=navy,clip=(-4:4,-4:4)] (\lc);
  \segment[draw=navy,clip=(-4:4,-4:4)] (\mc);

  \foreach \p/\placement in {A/above right,B/above left,C/above,
  D1/below right,E1/above,F1/below left,
  D2/above left,E2/above right,F2/below left}{%,
  %P1/above,P2/above,Q1/above,Q2/above,R1/above,R2/above}{
   \fill[red] (\p) circle (1.5pt);
   \draw (\p) node[\placement] {$\p$};
  }
\end{tikzpicture}

\end{document}